\scanpage{7}
% archive foliation: 6
D'autre part, la donnée de $(\mathcal X,\mathcal X_0)$ équivaut à celle de
\[
(\mathcal X_0,\mathcal X_1,\varphi)\qquad
\overset{\substack{\inserted{[f_*]}\\\Vert}}{\varphi}:\mathcal F(\mathcal X_0)\longrightarrow\mathcal F(\mathcal X_1)
\]
foncteur exact à g. \inserted{``accessible''}

On \unsure{s'intéresse|s'intéressera} à des conditions \inserted{sur $(\mathcal X_0,\mathcal X_1,\varphi)$} qui assurent que les \inserted{deux} inclusions
\diagram{\begin{tikzcd}
\mathcal X_0 \arrow[r] & \mathcal X & \mathcal X_1 \arrow[l]
\end{tikzcd}}
sont des équivalences faibles. Les cas qui nous intéressent le plus, sont les cas où $\varphi\inserted{=f_*}$ commute en fait aux $\underset{\leftarrow}{\lim}$ quelconques \inserted{i.e. a un adjoint à droite $f^*$}, --- comme p.ex. celui du § précédent, où on supposait \uline{de plus} que $f^*$ est exact à gauche (mais on \uline{n'a} \uline{pas} envie de le supposer ici). Un cas particulièrement important pour nous est celui où $\mathcal X$ est de la forme $\deleted{\unsure{\widehat X}}\inserted{\operatorname{Top}(X)}$, un $X$ dans $(\operatorname{Cat})$, il revient au $=$ de dire que
